\paragraph{Code LLMs and their benchmarks.}
Code-generation models are usually measured by functional correctness on small, self-contained
problems---HumanEval and its successors \cite{chen2021codex}---or, more recently, by repository-level
issue resolution \cite{jimenez2024swebench}. Neither rewards performance: a program that passes the
tests scores the same whether it runs in one second or ten. Our setting inverts this. Correctness is a
gate rather than the objective, and the score is measured wall-clock time of a real application.

\paragraph{LLMs for performance optimization.}
Performance-oriented work with LLMs falls into three groups. The first learns edits that make code
faster: PIE mines human performance-improving edits from competitive-programming submissions and adapts
LLMs to reproduce them, reporting large speedups on simulated CPU time \cite{shypula2024pie}. The
second targets the compiler rather than the source: Cummins et al.\ train models to predict
optimization flag sequences over LLVM-IR \cite{cummins2023llmopt} and release compiler foundation
models \cite{cummins2024llmcompiler}. The third targets GPU kernels directly. KernelBench asks models
to write CUDA kernels equivalent to PyTorch reference modules and reports that frontier models beat the
PyTorch baseline in fewer than 20\% of tasks \cite{ouyang2025kernelbench}; ParEval evaluates parallel
code generation across execution models and shows a marked gap between sequential and parallel
generation quality \cite{nichols2024parallel}. All of these generate or rewrite \emph{standalone}
kernels or short programs. The distinguishing feature of our study is that the artifact under
optimization is one file inside a production application---AutoDock-GPU, LAMMPS, Cholla---which keeps
its build system, its callers, and its own scientific validation.

\paragraph{Profiling data as LLM context.}
The closest work to ours feeds profiler output back to the model. Gai et al.\ route Nsight Compute,
SASS and Nsight Systems data through a roofline-based classifier that turns raw counters into natural
language guidance, and report large gains on KernelBench \cite{gai2026microprofiling}. We ask a
narrower and, we argue, prior question: does profiling data help \emph{at all}, relative to the same
prompt without it? The distinction matters most against KernelPro, the closest work to ours, which
feeds Nsight Compute, SASS and Nsight Systems output through pluggable micro-profiling tools that
convert raw metrics into natural-language guidance and reports a $125\%$ higher speedup for the
interpreted arm \cite{gai2026kernelpro}. Both of their arms carry profiling data; the comparison is
\emph{interpreted} against \emph{raw}. Ours removes the data entirely. Section~\ref{sec:guidvar}
evaluates their question too, on our benchmarks, and finds that interpretation widens the distribution
of outcomes rather than shifting its centre. Section~\ref{sec:profabl} isolates that by holding every other prompt component
fixed, and finds the answer depends on the workload---no measurable effect on single-kernel benchmarks,
but a clear effect on a large application where the bottleneck is not apparent from the source.

\paragraph{Autotuning and search.}
Classical autotuning searches a space of program variants or parameters with measurement in the loop:
OpenTuner for general program autotuning \cite{ansel2014opentuner}, and learned cost models for tensor
programs in AutoTVM \cite{chen2018autotvm} and Ansor \cite{zheng2020ansor}. These systems explore a
fixed, hand-designed space and need many measurements; an LLM instead proposes a small number of
semantically arbitrary source edits. The two are complementary, and our build--verify--measure loop is
essentially the autotuner's evaluation harness wrapped around a generative proposer. The profile
summaries we place in the prompt play the role that a cost model plays in search: they tell the
proposer where the time goes. Roofline analysis \cite{williams2009roofline} underlies both the
Nsight Compute reports we summarize and the bottleneck classification used by prior work.

\paragraph{Repair loops.}
Iterative self-repair from compiler and test feedback is now standard practice
\cite{chen2024selfdebug}. Our loop is a constrained instance: the feedback is a build log or an
application-level correctness gate, and the number of repair iterations is bounded. Our contribution
here is negative and practical---we document cases where the loop cannot work because the information
needed is absent from the feedback, for example a linker error whose cause line is filtered out of the
summary, or an application that aborts with a bare assertion and discards the CUDA error.

\paragraph{Benchmark gaming and reward hacking.}
A model optimizing against a checkable objective may satisfy the check without doing the work. This is
the code-generation instance of reward hacking, studied recently through benchmarks that deliberately
expose exploitable evaluation harnesses: EvilGenie places coding agents in an environment where
test files can be edited or answers hardcoded, and finds explicit reward hacking by current
proprietary agents, detectable by an LLM judge but not reliably by held-out tests
\cite{gabor2025evilgenie}. In the performance setting specifically, KernelBench-Verified audits
LLM-generated GPU kernels and finds identity functions reported as $374\times$ speedups, fast paths
gated on the test tensor's exact dimensions, and value-dependent work skipping; recalibrated against
hidden input distributions and an LLM-judge audit, the best model drops from $1.43\times$ to
$0.88\times$, with $11.6\%$ of problems relying on distribution-specific exploits
\cite{zhang2026kernelbenchverified}. Our contribution here is therefore not the observation but its
setting and its remedy: the same failure in an HPC benchmark rather than an ML one, a prompt-level
constraint that we measure, and a detection rule cheap enough to run on every candidate. We observe the phenomenon in a performance setting, where it
is particularly insidious: deleting the measured computation preserves the printed output and the
benchmark's own PASS check while producing three-order-of-magnitude ``speedups''
(Section~\ref{sec:intent}). We show that a prompt-level semantic-preservation constraint suppresses it
for both models we tested, and we report the detection rule we use in its place---an implausibility
threshold followed by diff inspection.

\paragraph{Repository-level optimization.}
A parallel line of work optimizes whole repositories rather than kernels. GSO scores agents against the
speedup an expert commit achieved across ten codebases \cite{shetty2025gso}, and PerfAgent adds a
profiler-guided, verifier-in-the-loop layer that more than doubles the rate of expert-matching patches
on it \cite{deng2026perfagent}. Its motivation matches ours almost exactly---agents miss bottlenecks
hidden behind abstraction layers and native extensions---but the targets are CPU-bound Python
repositories scored against an expert's commit. Ours are GPU kernels inside production scientific
applications, scored against the application's own scientific output, which is a correctness standard
an expert diff cannot provide.

\paragraph{HPC benchmark suites.}
We use HeCBench \cite{jin2023hecbench} for the prompt ablation, which provides many small CUDA kernels
with built-in verification, and three production applications for the end-to-end study. The contrast
between the two is one of our main findings: conclusions drawn on single-kernel benchmarks did not
transfer to applications, either for the value of profiling data or for the prompt variant that works
best.
